\documentclass[11pt,a4paper]{article}

% ============================================================
% Uniform MTDF preamble -- identical across all papers
% ============================================================
\usepackage[utf8]{inputenc}
\usepackage[T1]{fontenc}
\usepackage{lmodern}
\usepackage[english]{babel}
\usepackage[a4paper,margin=2.2cm]{geometry}
\usepackage{amsmath,amssymb,amsthm}
\usepackage{graphicx}
\usepackage{float}
\usepackage{booktabs}
\usepackage{array}
\usepackage{longtable}
\usepackage{tabularx}
\usepackage{xltabular}
\usepackage{multirow}
\usepackage{caption}
\usepackage[dvipsnames,table]{xcolor}
\usepackage{mdframed}
\usepackage{parskip}
\usepackage{ragged2e}
\usepackage[hidelinks,breaklinks=true]{hyperref}
\usepackage{microtype}
\usepackage{url}
% Allow \path{} to break more aggressively inside narrow table columns:
\makeatletter
\g@addto@macro\UrlBreaks{\do\_\do\-\do\.\do\/\do\:}
\makeatother

% Colours matching the HTML theme
\definecolor{accent}{HTML}{1A4B8A}
\definecolor{callout-bg}{HTML}{FBFBFB}
\definecolor{tablehead}{HTML}{F0F0F0}

% Prevent any table from overflowing
\setlength{\tabcolsep}{4pt}
\renewcommand{\arraystretch}{1.15}

% Caption style (compact, italic, small like the HTML)
\captionsetup{font=small,labelfont=bf,labelsep=period,justification=centerfirst}

% Hyperref metadata placeholders (filled per-paper below)
\hypersetup{
  pdftitle={MTDF Validation Suite Appendix: Reproducibility Index and Test Outcomes},
  pdfauthor={Ingo Mesche},
  pdfsubject={MTDF - Mesche Time-Dilation Field Theory},
  colorlinks=false
}

% Prevent overfull hbox from long URLs/DOIs in bibliography
\sloppy
\emergencystretch=3em

% Tolerant line breaking so long identifiers (Python filenames etc.) don't
% produce large overfull hboxes in narrow table columns. These settings keep
% all content on-page while slightly relaxing right-margin strictness.
\tolerance=1200
\hbadness=2000
\hyphenpenalty=150
\exhyphenpenalty=50

% ============================================================
\title{MTDF Validation Suite Appendix: Reproducibility Index and Test Outcomes}
\author{Ingo Mesche\\ \small Independent Researcher, Malta \\ \small \texttt{imesche@outlook.com}}
\date{V75 / July 2026}

\begin{document}
\maketitle
\begin{center}\itshape Shared appendix referenced by MTDF Companion Papers Parts I to IV. Updated for Validation\_\allowbreak{}Dashboard\_\allowbreak{}V75.html.\end{center}




\begin{abstract}
This appendix consolidates the reproducibility index, test scripts, and summary diagnostics for the MTDF validation suite. It provides the complete inventory of tests referenced across the companion papers (MTDF\_\allowbreak{}01 through MTDF\_\allowbreak{}08), including script locations, input data, and pass/\allowbreak{}fail outcomes under the V18 strict protocol.
\end{abstract}

\begin{mdframed}[linecolor=accent,linewidth=0.5pt,backgroundcolor=callout-bg,leftline=true,rightline=true,topline=true,bottomline=true,innerleftmargin=8pt,innerrightmargin=8pt,innertopmargin=6pt,innerbottommargin=6pt]
\textbf{Claim status.}
This paper distinguishes between:
(i) calibration anchors,
(ii) validation targets,
(iii) diagnostic consistency checks, and
(iv) speculative extensions.
Only results explicitly labelled as validation targets are used as evidence for the core MTDF claim. Diagnostic and speculative sections are included to define future tests, not to claim established confirmation.

\textbf{Conventions.} Growth and S8 material defers to the sealed v2 chains (the v1 Phase-5 report is retired as a gauge artifact; see MTDF\_\allowbreak{}01 Section 7.3 and MTDF\_\allowbreak{}07 Section 6).
\end{mdframed}


\section{Purpose and how to use this appendix}
This appendix consolidates the reproducibility index and quick-look outcomes for the test scripts used across the MTDF companion papers.
  The intent is simple: keep the core papers readable, while keeping the test inventory, scripts, and summary diagnostics in one stable place.

Files referenced here:

\begin{itemize}
  \item Companion Part I: Environmental and local evidence
  \item Companion Part II: Gravity sector and lensing validation (\textit{MTDF\_\allowbreak{}03}~[see companion paper])
  \item Companion Part III: Photon coupling, redshift, and early universe links
  \item Companion Part IV: Cosmological validation and high-energy test programme
  \item Interactive dashboard: {\small\protect\path{Validation_Dashboard_V75.html}}
\end{itemize}

\begin{mdframed}[linecolor=accent,linewidth=0.5pt,backgroundcolor=callout-bg,leftline=true,rightline=true,topline=true,bottomline=true,innerleftmargin=8pt,innerrightmargin=8pt,innertopmargin=6pt,innerbottommargin=6pt]
\textbf{Scope note:}
    This appendix is a reproducibility index. It does not change the evidential status of any result. Claim status is defined in the corresponding main or companion paper. This appendix records scripts and diagnostics. Statistical weight remains with the distance-calibrated observables and the calibration versus validation separation defined in the dashboard workflow.
\end{mdframed}

\section{Test index and reproducibility}
\label{sec:appendix-a}
This appendix records the scripts used to generate the figures and summary diagnostics for Parts I to IV of the broader MTDF validation effort.
The SN Ia $\times$ void result is primary evidence within Companion Part I's extension programme; it is not part of the V18 strict validation core. The remaining tests are included as constraints, consistency checks, and forward-looking discriminants across Parts II and III.

\subsection{Script index by paper component}
\begingroup\scriptsize
\begin{xltabular}{\linewidth}{>{\RaggedRight\arraybackslash}X>{\RaggedRight\arraybackslash}X>{\RaggedRight\arraybackslash}X>{\RaggedRight\arraybackslash}X}
\toprule
\textbf{Paper component} & \textbf{Test} & \textbf{Script} & \textbf{Role} \tabularnewline
\midrule\endfirsthead
\toprule
\textbf{Paper component} & \textbf{Test} & \textbf{Script} & \textbf{Role} \tabularnewline
\midrule\endhead
\bottomrule\endfoot
Part I (environment and independent distances) & SN Ia $\times$ void analysis & sn\_\allowbreak{}void\_\allowbreak{}*.py & Primary evidence for an environment-dependent distance-calibrated residual. \tabularnewline
\multirow{8}{=}{Part III (photon coupling $\kappa$) \emph{[speculative]}} & Zero-velocity shell & {\small\protect\path{mtdf_zero_velocity_shell_test.py}} & Null and pipeline integrity check. \tabularnewline
 & Lensed image $\Delta$z & {\small\protect\path{mtdf_lensing_redshift_split_test.py}} & Path-dependence test using multiple lensed images of the same source. \tabularnewline
 & Achromaticity & {\small\protect\path{mtdf_achromaticity_test.py}} & Wavelength independence requirement for any stress-induced redshift. \tabularnewline
 & Dipole mismatch & {\small\protect\path{mtdf_dipole_mismatch_test.py}} & Checks whether a stress-induced dipole can contribute to the quasar dipole anomaly. \tabularnewline
 & Dipole alignment & {\small\protect\path{mtdf_dipole_vector_alignment.py}} & Directional consistency check between excess dipole and density gradient. \tabularnewline
 & Distance duality & {\small\protect\path{mtdf_distance_duality_test.py}} & Etherington relation and potential environment split. \tabularnewline
 & Sandage-Loeb drift & {\small\protect\path{mtdf_sandage_loeb_drift_test.py}} & Long-baseline, future definitive test of evolving redshift physics. \tabularnewline
 & Synthesis plot & {\small\protect\path{mtdf_synthesis_plot.py}} & Conceptual summary figure linking multiple anomalies to the same scaling law. \tabularnewline
\multirow{6}{=}{Part IV (cosmological and lensing)} & S8 tension and growth & {\small\protect\path{mtdf_growth_S8_analysis.py}} & Large-scale structure growth consistency. \tabularnewline
 & BAO void split & {\small\protect\path{mtdf_bao_void_split_test.py}} & Environment split for BAO-scale observables. \tabularnewline
 & GW standard sirens & {\small\protect\path{mtdf_gw_siren_test.py}} & Independent distance ladder cross-check using gravitational waves. \tabularnewline
 & Lensing versus RSD (E\textsubscript{G}) & {\small\protect\path{mtdf_lensing_rsd_test.py}} & Modified-gravity consistency test, not a direct probe of photon coupling. \tabularnewline
 & Cluster mass MDR & {\small\protect\path{mtdf_cluster_mass_discrepancy_test.py}} & Compares dynamical and lensing mass proxies in clusters. \tabularnewline
 & AP void shape & {\small\protect\path{mtdf_ap_void_test.py}} & Alcock-Paczynski void geometry constraint. \tabularnewline
\multirow{3}{=}{Part II (gravity sector and lensing)} & Strong lensing time delays & {\small\protect\path{step20_strong_lensing.py}} & Time-delay distance consistency check under eta = 1. \tabularnewline
 & Elliptical galaxy dispersions & {\small\protect\path{step22_elliptical_dispersions.py}} & Pressure-supported system check (SLACS-style). \tabularnewline
 & Satellite kinematics & {\small\protect\path{step22b_satellite_kinematics.py}} & Digitised and table-based comparison (More+2011 style) with robustness notes. \tabularnewline
\end{xltabular}
\endgroup

\subsection{Outcome snapshot}
\begingroup\scriptsize
\begin{xltabular}{\linewidth}{>{\RaggedRight\arraybackslash}X>{\RaggedRight\arraybackslash}X>{\RaggedRight\arraybackslash}X}
\toprule
\textbf{Script} & \textbf{Headline outcome} & \textbf{Interpretation at current sensitivity} \tabularnewline
\midrule\endfirsthead
\toprule
\textbf{Script} & \textbf{Headline outcome} & \textbf{Interpretation at current sensitivity} \tabularnewline
\midrule\endhead
\bottomrule\endfoot
{\small\protect\path{mtdf_zero_velocity_shell_test.py}} & Consistent & No significant artifact detected in a designed null scenario. \tabularnewline
{\small\protect\path{mtdf_distance_duality_test.py}} & Consistent & No significant violation of $\eta$ = 1 at current precision; environment split remains a future target. \tabularnewline
{\small\protect\path{mtdf_lensing_redshift_split_test.py}} & Consistent & No significant $\Delta$z between multiple lensed images; current spectroscopy is above the $\kappa$ $\lesssim$ 10\textsuperscript{-3} signal level. \tabularnewline
{\small\protect\path{mtdf_achromaticity_test.py}} & Consistent at current precision & No evidence for chromatic redshift; any measured offsets are likely systematics at the present 10\textsuperscript{-5} level. \tabularnewline
{\small\protect\path{mtdf_dipole_mismatch_test.py}} & Suggestive (partial amplitude) & At $\kappa$ $\approx$ 0.001, MTDF reproduces the correct sign and roughly 42\% of the reported quasar dipole excess amplitude. \tabularnewline
{\small\protect\path{mtdf_dipole_vector_alignment.py}} & Moderate directional agreement & Alignment is not tight enough to be a strong constraint, but it remains informative for future catalogues. \tabularnewline
{\small\protect\path{mtdf_sandage_loeb_drift_test.py}} & Future-only & Predicted environment split in \.{z} is far below current detection limits for the derived $\kappa$ value (f\textsubscript{kick}/3). \tabularnewline
{\small\protect\path{mtdf_synthesis_plot.py}} & Key synthesis figure & Heuristic summary connecting multiple anomalies across environments. It is not, by itself, a measurement of $\kappa$. \tabularnewline
{\small\protect\path{mtdf_lensing_rsd_test.py}} & Consistent & E\textsubscript{G} data are broadly consistent with GR at current precision; not a direct $\kappa$ probe. \tabularnewline
{\small\protect\path{mtdf_cluster_mass_discrepancy_test.py}} & Historical exploratory row; under the frozen architecture $\eta$ = 1 gives M\textsubscript{dyn}/M\textsubscript{lens} = 1.000 exactly (see MTDF\_\allowbreak{}07, Test C5): this channel is not a discriminator. & Historical exploratory note; superseded by the $\eta$ = 1 result (MTDF\_\allowbreak{}07, Test C5). \tabularnewline
{\small\protect\path{mtdf_ap_void_test.py}} & Consistent & Void shape residuals consistent with zero after RSD correction; $\kappa$ effects are below current precision. \tabularnewline
{\small\protect\path{mtdf_growth_S8_analysis.py}} & See Part IV (Phase 5) & Growth and S8 are treated in Part IV via the sealed v2 full-Planck chains (the v1 Phase-5 report is retired as a gauge artifact); this script is an earlier exploratory check and its results are not used there. \tabularnewline
{\small\protect\path{mtdf_bao_void_split_test.py}} & Exploratory, not in the papers & Exploratory diagnostic; BAO environment splits are not part of the validated core and are not presented in Part IV. \tabularnewline
{\small\protect\path{mtdf_gw_siren_test.py}} & Exploratory, not in the papers & Exploratory diagnostic; GW standard-siren distance checks are not presented in the papers and carry no evidential weight here. \tabularnewline
{\small\protect\path{step20_strong_lensing.py}} & See Part II & Time-delay distance chi\^{}2/\allowbreak{}nu = 0.77; isothermal slope gamma = 2.000 (structural prediction). See gravity-sector companion. \tabularnewline
{\small\protect\path{step22_elliptical_dispersions.py}} & See Part II & Mean velocity dispersion bias -1.4\%; chi\^{}2/\allowbreak{}nu = 1.81 (gas-corrected). See gravity-sector companion. \tabularnewline
{\small\protect\path{step22b_satellite_kinematics.py}} & See Part II & More+2011 amplitude test: mean |bias| 6.3\%; profile chi\^{}2/\allowbreak{}nu = 3.60 (best nuisance). See gravity-sector companion. \tabularnewline
\end{xltabular}
\endgroup

\subsection{Synthesis figure}
Figure A.1 provides a compact visual summary that places three anomalies on a single stress-scaling trend.
This figure is intended as a high-level synthesis and motivates a unified search strategy. It should not be interpreted as a standalone parameter inference for $\kappa$.

\begin{figure}[H]
\centering
\includegraphics[width=\linewidth,keepaspectratio]{figures/mtdf_06_figure_a_1_the_stress_redshift_main_sequ.png}
\caption{\textbf{Figure A.1.} The stress-redshift main sequence synthesis plot (generated by {\small\protect\path{mtdf_synthesis_plot.py}}).}
\end{figure}

\subsection{P1-SPARC test definition}
\textbf{Statistic.} The P1-SPARC benchmark reports the aggregate root-mean-square of the log-space residual log\textsubscript{10}(v\textsubscript{obs}/v\textsubscript{model}) over the full SPARC sample (Lelli et al. 2016, AJ 152, 157; 175 galaxies). The achieved value is 0.185 dex against a target of 0.174 $\pm$ 0.011 dex, below the 0.25 dex reference threshold. This is an aggregate amplitude statistic, not a per-galaxy rotation-curve shape fit, and no per-galaxy $\chi$$^{2}$ of zero is claimed.

\textbf{Model and parameters.} The MTDF velocity is v\textsubscript{model}(r) = v\textsubscript{bar}(r) $\sqrt{\,}$(1 + $\alpha$/(1 + r/\allowbreak{}$\beta$)) with $\alpha$ = 1.30 and $\beta$ = 22.7 Mpc frozen. Because $\beta$ greatly exceeds galactic radii, the enhancement is effectively a constant factor $\sqrt{\,}$(1+$\alpha$) $\approx$ 1.52 in velocity (a factor 1+$\alpha$ $\approx$ 2.3 in acceleration) across the rotation curve.

\textbf{Mass-to-light treatment and error model.} v\textsubscript{bar} is built from the SPARC mass-model decomposition (gas, disk, bulge); the aggregate scatter is evaluated at the mass-model normalisation, and the residual uses the observed velocity uncertainties. For reference, a per-galaxy shape fit at the fixed SPARC-standard mass-to-light ratios (0.5 for disks, 0.7 for bulges at 3.6 $\mu$m) yields per-galaxy reduced chi-squared of order 10--280 (e.g. NGC~3198 $\approx$ 72, NGC~2403 $\approx$ 57, NGC~7331 $\approx$ 281), confirming that the constant enhancement reproduces the amplitude of the rotation curves but not their shape.

\textbf{Shape diagnostic (radial acceleration relation).} On the radial acceleration relation, the observed data follow a low-acceleration log-log slope of $\approx$ 0.58 (consistent with the McGaugh et al. 2016 value near 0.5), whereas a constant enhancement predicts a slope of 1.0 (a parallel offset). MTDF therefore matches the approximate amplitude of the relation (median g\textsubscript{obs}/g\textsubscript{bar} $\approx$ 2.7 against a predicted 1 + $\alpha$ = 2.3) but not its curvature; the constant-enhancement parameterisation used here does not bend. The curvature is produced by the low-Q branch of the frozen unified action (MTDF\_\allowbreak{}01, Section 2A), which resolves the shape-level question; this appendix's parameterisation is retained as the historical effective description.

\textbf{Scope.} The galaxy-scale pillars are treated as scalar benchmarks and are excluded from the V18 strict scoring totals, so the shape-level limitation above does not affect the strict vector results. Audit artefacts (per-galaxy table and the radial-acceleration-relation diagnostic) are provided in the reproducibility package under \texttt{validation/\allowbreak{}output/\allowbreak{}audit\_\allowbreak{}sparc/\allowbreak{}}.

\section{Data and code availability}
The complete MTDF reproducibility package, including the modified CLASS solver
(\texttt{class\_\allowbreak{}mtdf}), the V19 strict validation workbook, the
\texttt{{\small\protect\path{run_validate.py}}} pipeline, the gravity-sector artefact manifest,
and the Phase 2 CosmoPower emulator cache, is archived on Zenodo
(concept DOI \href{https://doi.org/10.5281/zenodo.19741058}{{\small\protect\path{10.5281/zenodo.19741058}}};
this release, \texttt{v1.1.7}, DOI
\href{https://doi.org/10.5281/zenodo.21525855}{{\small\protect\path{10.5281/zenodo.21525855}}})
and mirrored on GitHub at
\href{https://github.com/IMesche/MTDF}{github.com/\allowbreak{}IMesche/\allowbreak{}MTDF}.
The \texttt{v1.1.7} Git tag corresponds to the archived Zenodo release. All
numerical results reported here can be regenerated from the shipped workbook
and scripts following the procedure in the top-level \texttt{{\small\protect\path{README.md}}}.


\end{document}
